\section{Formal Substrate-Gap Propositions}
\label{app:formal-substrate-gap}

This appendix expands the formal statements in
Section~\ref{sec:alignment}.  The statements are intentionally narrow.  They
support the trace-bound diagnostic framing and the spoofability interpretation
of the tested public predicates; they do not prove a universal impossibility
theorem for all future public final-text predicates.

\subsection{First-Divergence Reduction}

\begin{proof}[Proof of Proposition~\ref{prop:first-divergence-reduction}]
Let $C$ be the deterministic evidence rule, and let $u$ and $v$ be tokenized
continuations such that $C(u)\neq C(v)$.  Since a deterministic rule assigns
the same class to identical inputs, $u$ and $v$ cannot be identical token
sequences.  By Lemma~\ref{lem:first-divergence}, distinct token sequences have
a first-divergence position after their longest shared prefix.  The two
branches at that position are therefore a witness that the continuations have
already become distinguishable at tokenizer-native next-token resolution.
\end{proof}

The proposition does not say that the first divergent token alone is sufficient
to decode a payload or verify a claim.  It only states that any deterministic
class distinction between two tokenized continuations must begin at some
next-token branch.  This is why first-divergence traces are useful as an
upper-bound diagnostic for provider-side event controllability.

\subsection{Source-Mismatch Accept Density}

\begin{proof}[Proof of Proposition~\ref{prop:public-predicate-spoofability}]
Draw independent samples $x_1,x_2,\ldots$ from the source-mismatch proposal
distribution $D$ and evaluate the deterministic public predicate $A$ on each
sample.  Each trial succeeds with probability
$p=\Pr_{x\sim D}[A(x)=1]$.  The number of trials until the first accepted sample
is therefore a geometric random variable with expectation $1/p$ when $p>0$.
The accepted sample came from $D$, a source that should not satisfy the
protected claim.  Its acceptance is therefore evidence that the predicate has a
source-mismatch accept region.  It is not protected success, not codeword
recovery, and not a trace-bound recovery result.
\end{proof}

The result is deliberately conditional on the tested predicate and proposal
distribution.  It does not rule out every possible future public predicate.
It explains why the predicates evaluated in this manuscript must be treated as
spoofing surfaces when they have nonzero source-mismatch accept density.

\subsection{Substrate Displacement Scope}

Principle~\ref{prin:substrate-displacement} is a claim-scope rule.  A
metadata manifest, protocol transcript, proof object, provider trace, or
model-state witness can carry stronger evidence than ordinary final text.
However, the stronger result applies only when the corresponding substrate is
available to the verifier.  A copied-text dispute that lacks that substrate
cannot inherit the same deterministic witness merely because the stronger
substrate exists elsewhere.

This principle is not a no-go theorem for all verification systems.  It is a
discipline for reading the experiments in this paper: trace-bound recovery is
evidence of provider-side event controllability, while public final-text
predicate accepts from source-mismatch examples are evidence of spoofability in
the tested public-predicate substrate.
